%=======================================================================
% Requires: booktabs, multirow, amsmath, amssymb, siunitx
%=======================================================================
\section{Results}
\label{sec:results}

The pipeline of Section~\ref{sec:collection} is applied to sequence
\texttt{coop2\_0828} under both ego assignments. Every quantity below is
produced by the released scripts from the recording, the anchored map, and the
annotation file; none is asserted from the collection plan.

\subsection{Sequence Characterization}

Table~\ref{tab:seq} gives the descriptor defined in
Section~\ref{sec:descriptor}, one row per ego assignment, with the design fields
spanning both rows and the measured geometry reported separately for each.
Occlusion episodes have median duration \SI{18.8}{\second}, and the elevation
offset to the infrastructure node is \SI{1.53}{\metre} (IQR
\SIrange{1.52}{1.54}{\metre}) under both assignments.

\begin{table*}[t]
\centering
\caption{Per-sequence descriptor for \texttt{coop2\_0828}, both ego assignments.
Columns are as defined in Section~\ref{sec:descriptor}. $^\ast$ staged ego.
Episodes are contiguous intervals of $\neg\,\mathrm{visible}(e)$ of at least
\SI{0.5}{\second}, with the number covered by a helper in parentheses. $\beta$
is measured to the infrastructure node. Range rescue is reported under RGB-D
evidence. NLoS aggregates obstructed and degraded frames on the ego--AP link.
An em-dash denotes a quantity undefined for the sequence, not a zero.}
\label{tab:seq}
\small
\setlength{\tabcolsep}{4pt}
\begin{tabular}{@{}lllllccccccccc@{}}
\toprule
\multicolumn{5}{c}{Design} & \multicolumn{4}{c}{Geometry} &
\multicolumn{3}{c}{Rescue rate} & \multicolumn{2}{c}{Link, scene} \\
\cmidrule(r){1-5}\cmidrule(lr){6-9}\cmidrule(lr){10-12}\cmidrule(l){13-14}
Sequence & Site & Paradigm & Stress & Ego &
$C$ & $\beta$ ($^\circ$) & $\beta > 90^\circ$ & Episodes &
Occl. & Range & Dopp. &
NLoS (\%) & Obj. \\
\midrule
\multirow{2}{*}{\texttt{coop2\_0828}} & \multirow{2}{*}{MIRC} &
\multirow{2}{*}{co-observation} & \multirow{2}{*}{---} &
\texttt{m1}$^\ast$ & 0.5 & 149 & 0.67 & 9 (6) & 0.10 & 0.29 & --- & 33 & 2 \\
 & & & & \texttt{m2} & 0.5 & 29 & 0.33 & 8 (5) & 0.02 & 0.39 & --- & 18 & 2 \\
\bottomrule
\end{tabular}
\end{table*}

The run was planned as co-observation. Measured, it is co-observation in
74\,\% of ego-in-FoV observations with \texttt{mobile\_1} as ego and 77\,\% with
\texttt{mobile\_2}; the remainder are occlusion, handover, and unstructured
frames that the plan did not call for and that the labels record regardless.
This is the design check of Section~\ref{sec:paradigms}: a trajectory written
for one condition delivered that condition in approximately three frames out of
four, and the release identifies which frames were the fourth.

Doppler diversity is left empty rather than reported as zero. The predicate is
undefined in the absence of moving targets, and the sequence contains none.

\subsection{Validation of the Released Labels}

The occlusion and link-state labels are geometric predictions derived from a
map. Each is therefore checked against a measured signal that did not enter its
computation (Table~\ref{tab:validation}).

\begin{table}[t]
\centering
\caption{Validation of the geometric labels against independent measurement;
$\rho$ is the Spearman rank correlation over in-FoV observations.}
\label{tab:validation}
\small
\begin{tabular}{@{}lll@{}}
\toprule
Label & Independent check & Result \\
\midrule
Occlusion (\texttt{m1}) & Ouster / ZED points in box & $\rho = -0.46$ / $-0.44$ \\
Occlusion (\texttt{m2}) & RealSense depth points in box & $\rho = -0.13$\,$^\dagger$ \\
Link clearance (\texttt{m1}) & measured RSSI & $\rho = +0.73$; LOS $-35$ vs NLOS $-43$\,dBm \\
Link clearance (\texttt{m2}) & measured RSSI & $\rho = +0.67$; LOS $-39$ vs NLOS $-50$\,dBm \\
\bottomrule
\end{tabular}
\end{table}

Predicted occlusion is negatively correlated with the evidence actually
returned, on both of \texttt{mobile\_1}'s independent modalities and at
comparable strength, which is the expected signature of a correct occlusion
label against noisy sensors. Predicted link clearance is strongly positively
correlated with measured RSSI on both mobile nodes, and the clearance-derived
LOS/NLOS partition separates median RSSI by \SI{8}{\decibel} and
\SI{11}{\decibel} respectively: a purely geometric predicate reproducing a radio
measurement that did not enter it.

\noindent$^\dagger$ Uninformative rather than negative. \texttt{mobile\_2}
exceeded $\tau$ in 2\,\% of its observations, leaving almost no variation in the
predicted quantity to correlate against. Its evidence chain was therefore
validated along the axis that does vary for that node, range, and the expected
collapse confirmed: \num{861} depth points within the box at \SI{2.4}{\metre},
none beyond \SI{7}{\metre}. The row indicates that the check could not be run,
not that the label failed.

\subsection{Findings}

\paragraph{Range, not occlusion, is the binding constraint.}
The condition the collection was designed around is not the condition that
limits the team. Occlusion rescue, in which the ego is blocked and a helper is
not, occurs on 0.10 and 0.02 of ego-in-FoV observations. Range rescue, in which
the ego holds a clear line but returns too little evidence while a partner does
not, occurs on 0.29 and 0.39. Stratified by ego--object range, rescue rates are
near zero below \SI{5}{\metre} and rise to \numrange{0.5}{0.7} beyond it. The
mechanism appears directly in \texttt{mobile\_2}'s per-range statistics: across
the \SIrange{2}{5}{\metre}, \SIrange{5}{10}{\metre}, and
\SIrange{10}{15}{\metre} bins its fraction \textit{visible} remains at
\numrange{0.96}{1.00} while its fraction \textit{observed} collapses
\num{0.85} $\rightarrow$ \num{0.00} $\rightarrow$ \num{0.00}. The geometry
admits sight of the object throughout; the sensor ceases to return it past the
first bin.

This is the separation of~\eqref{eq:visible} from~\eqref{eq:observed} doing
work. A release reporting line-of-sight geometry alone would characterise this
team as having near-complete coverage. The effect was not staged, which makes it
stronger evidence than the occlusion episodes that were.

\paragraph{The choice of evidence source changes the rate by a factor of 36.}
\texttt{mobile\_1}'s range-rescue rate is \num{0.008} when its evidence is
Ouster LiDAR returns and \num{0.29} when it is ZED stereo depth, for the same
trajectory, scene, ray casts, annotations, and occlusion labels. Only the
definition of sufficient evidence changed, through the modality's floor and its
calibrated range cap. A collaboration rate reported without its evidence source
is therefore not reproducible to better than a factor of tens, and rates
computed under different sensing budgets are not comparable. The per-modality
evidence counts are released per (frame, node, object) so that any such rate may
be recomputed under a stated definition, and Table~\ref{tab:seq}'s range-rescue
column is reported under the RGB-D definition explicitly.

\paragraph{Viewpoint complementarity is positional, not elevational.}
The same infrastructure node, in the same pose, constitutes a complementary
viewpoint for one mobile node and a redundant one for the other. With
\texttt{mobile\_1} as ego the median aspect angle to infrastructure is
$149^\circ$, with $\beta > 90^\circ$ on 67\,\% of frames; with
\texttt{mobile\_2} as ego it is $29^\circ$ and 33\,\%. The elevation offset is
identical in both cases ($\Delta h = \SI{1.53}{\metre}$, IQR
\SIrange{1.52}{1.54}{\metre}); what differs is the ego's position relative to
the infrastructure and the object. Height confers persistence and an
unobstructed downward line, but not by itself the wide baseline from which
fusion benefits. Reporting $\Delta h$, as infrastructure-augmented datasets
conventionally do, does not characterise the collaboration geometry, whereas
$\beta$ conditioned on the ego does.
